% Non-overlap property — promote the methodological observation from
% docs/notes/materialisation-and-overlap.md into a paper section.
% Keep the salt warning here too — it documents which steps need to
% stay pure for the invariant to hold.

\emph{[TO BE REWORDED]}

The composition operators of Sections \ref{ssec:axis-flip} through
\ref{ssec:peel} are subject to the following structural invariant.

\begin{theorem}[Non-overlap]\label{thm:nonoverlap}
Let \(O\) be an outer bilinear scheme of shape
\(\nmpshape{n_o}{m_o}{p_o}\) with \(r_o\) products, and let
\(L_1, \ldots, L_{r_o}\) be sub-schemes assigned via
\code{Recombination.constructWithAllocation} to a target shape
\(\nmpshape{n}{m}{p}\) with allocations \((\mathbf{a}_A, \mathbf{a}_B,
\mathbf{a}_C)\) and an optional output peel pattern. The resulting
algorithm has rank exactly \(\sum_k r_{L_k}\), and no two of its
bilinear products are equal as bilinear forms.
\end{theorem}

\begin{proof}[Sketch]
The construction is block-substitution: each outer product
\(k \in [1, r_o]\) is instantiated as the sub-scheme \(L_k\) wired
into block-positions of \((A, B, C)\) determined by the allocations
and the outer scheme's bilinear form. Two columns of the resulting
\((U, V, W)\) can be equal as bilinear forms only if (a) they come
from the same outer product index AND (b) they correspond to the
same column of \(L_k\) -- which would require the same column of the
sub-scheme to be wired to two different positions, contradicting the
definition of block-substitution. The peel reduction
(Section \ref{ssec:peel}) removes dead rows of \(W\) but never
identifies two columns. The axis-flip orbit (Section
\ref{ssec:axis-flip}) is an outer transformation that does not change
the rank or the column-distinctness.
\end{proof}

\paragraph{Scope: the operations that \emph{reduce} below the sum.} The
theorem is about pure block-substitution, so it governs Kronecker,
concatenation, recombination-with-allocation, the axis-flip orbit and the
peel. Three operations deliberately fall \emph{outside} it by producing
fewer than \(\sum_k r_{L_k}\) products --- and none is a counterexample,
because none \emph{discovers} sharing; each applies a known,
\emph{predict-time} reduction.
\begin{itemize}[itemsep=2pt]
  \item \textbf{Pair fusion (Pan TA)} is not a composition operator at
    all: it is an optimisation on a decomposition's product set, fusing
    complementary pairs through Pan's trilinear-aggregation
    \emph{identity}. The saving is a property of that fixed identity,
    counted before materialisation --- TA does not break the invariant,
    it just is not a \(\sum_k r_{L_k}\) construction.
  \item \textbf{The serendipitous product} replaces an enlarged sub-block
    by its catalogued best rank --- again a lookup performed at predict
    time.
  \item \textbf{Downward projection} restricts a materialised parent and
    runs DCE; the resulting rank \(R-\mu\) is exactly what the
    projection-margin score (Section~\ref{sec:search}) \emph{predicts}
    from the parent's support. The DCE is the same deterministic step in
    both the prediction and the materialisation, so nothing is discovered.
\end{itemize}
So the deeper invariant below --- \emph{predicted rank equals materialised
rank, with no sharing found during materialisation} --- holds for these
three as well, even though the literal \(\sum_k r_{L_k}\) equality of the
theorem does not.

Theorem \ref{thm:nonoverlap} has two practical consequences worth
calling out.

\paragraph{Predicted rank equals materialised rank.} The frontier-closure
search (Section \ref{sec:search}) propagates rank predictions
\(\hat r = \sum_k r_{\text{sub-}k}\) computed by summing sub-shape
lookups in the catalog. The materialised scheme has exactly that
rank. There is no hidden discovery during materialisation; the
materialisation step is faithful execution of whatever strategy the
search picked. This is what lets us split the closure loop into a
fast rank-only search phase and a lazy materialise-only phase
without losing accuracy -- see Section \ref{sec:search} for the
algorithm.

\paragraph{A methodological boundary against hand-crafted schemes.}
Almost every hand-crafted small-format result in the literature
exists \emph{because} its author found a non-obvious shared
bilinear core: Strassen's \(\nmpfield{\Reals}{2}{2}{2} = 7\)
beats 8 because two of his seven products do double duty across
output cells; Laderman's \(\nmpfield{\Reals}{3}{3}{3} = 23\)
is the same story at larger scale; Smirnov, Pan, AlphaTensor,
AlphaEvolve all hunt for non-obvious cross-product sharing.
Theorem \ref{thm:nonoverlap} says our \emph{composition} operators do
\emph{not} hunt for such sharing -- the wins come from picking the
right composition of cataloged building blocks. (The reducing operators
of Section~\ref{ssec:decomposition} ff.\ --- serendipity, pair fusion,
projection --- do exploit known reductions, but never discover new ones;
see the scope remark above.)

This makes a ``composition matches hand-crafted rank'' outcome a
\emph{co-existence} claim, not a re-discovery claim. When the search
reaches \(\nmpfield{\Rationals}{17}{17}{17} = 2930\) via a Winograd
recombination at \((9,8)^3\), it does so by a route structurally
distinct from the hand-crafted disjoint-sum construction for the same
shape (FMM-Lille, rank \(2934\)): our scheme is a plain sum of
independent sub-products and reproduces none of that construction's
shared bilinear identities. The catalog records both with distinct
lineages and distinct \code{attribution\_for\_rank}.

The two approaches are complementary, not redundant. Composition
expands the SOTA frontier without finding new bilinear identities;
hand-crafted constructions find new bilinear identities that can be
ingested as new outer schemes, immediately expanding what
composition can do.

\paragraph{The invariant is current-state, not eternal.}
Theorem \ref{thm:nonoverlap} depends on \code{constructWithAllocation}
being pure block-substitution. Several extensions would break it:
post-construction column fusion that identifies duplicate bilinear
forms; random sign salting of sub-blocks that enriches the symmetry
orbit; coefficient mixing between outer products that share a
sub-shape. Any of these would let materialised rank be \emph{lower}
than predicted rank, which would be a feature -- but would also
break the search-rank-equals-materialised-rank invariant that the
two-phase search of Section \ref{sec:search} relies on. The
implementor of any such ``salt'' must either restrict it to the
final materialisation step (so the search keeps its pure prediction)
or re-derive predictions to account for the salt. The catalog
documents this in \code{docs/notes/materialisation-and-overlap.md}.

We have considered whether to add such salt and decided against it
for now: a pure search-time prediction is easier to reason about,
easier to cache, and avoids a class of subtle correctness bugs at
the search/materialise boundary. Hand-crafted-style sharing
discoveries are deferred to the one dedicated constructor we have ---
pair fusion (Section \ref{ssec:pair-fusion}) --- rather than emerging
from generic recombination; the disjoint-sum / $\tau$ route
(Section \ref{ssec:disjoint-sum}) remains future work.
